

```latex
\documentclass{article}
\usepackage{pathfinder-note}

\title{Competitive Market Behavior and Cooperative Survival: Cross-Paper Analysis}

\begin{document}

\section*{The Pair}

Two papers examine LLM agents in economic environments with contrasting outcomes. Q studies LLM agents in double auction markets, finding that agents exhibit slower or no convergence toward market equilibrium compared to humans, thus delivering less efficient allocations~\ledger{1}. P introduces The Energy Society, a survival economy where LLM agents cooperate under pressure, finding that cooperative incentives enable coordination, resource sharing, and collective problem-solving~\ledger{1}.

\section*{Connexion Pursued}

The hypothesis was to port P's cooperative mechanisms---cooperative objectives, memory systems, and action recommendations---into Q's double auction environment to test whether they restore convergence to equilibrium~\ledger{1}. The rationale: Q identifies a failure mode (LLM agents cannot discover equilibrium prices); P identifies success conditions (cooperation + memory + recommendations enable coordination).

\section*{Supported Findings}

The ledger establishes several findings about this proposed transfer:

\begin{itemize}
  \item Q's failure mode is strategic misalignment---agents cannot find equilibrium prices through bidding and asking behavior~\ledger{2}.
  \item P's memory system helps agents calibrate risk from past outcomes, which could address Q's observation that agents struggle with price adjustment strategy~\ledger{3}.
  \item P's action recommendations enable agents to suggest actions to peers; in Q's context, this could enable price signal sharing or strategy hints~\ledger{3}.
  \item The most promising transfer is memory plus recommendations without altering the auction's competitive payoff structure---testing whether information-sharing architecture alone improves convergence~\ledger{3}.
\end{itemize}

\section*{Objections}

A critical gap was identified: double auctions are inherently competitive mechanisms, while P's cooperation mechanisms assume a shared survival goal~\ledger{2}. The question remains whether cooperative incentives designed for collective survival can transfer to competitive price-discovery tasks, and this requires theoretical justification that the ledger does not provide~\ledger{2}. P's cooperative incentives---where agents care about group survival---conflict with the individual profit maximization inherent to double auctions~\ledger{3}.

\section*{Unresolved Issues}

The ledger does not resolve:
\begin{itemize}
  \item Whether information-sharing mechanisms (memory, recommendations) can improve convergence without changing payoff incentives.
  \item What theoretical framework would justify applying cooperative mechanisms to competitive environments.
  \item Whether heterogeneity across model families and market roles (observed in Q) would interact with memory or recommendation systems.
\end{itemize}

\section*{Smallest Next Step}

Implement memory and action recommendations in Q's agent architecture while keeping competitive payoffs unchanged. Measure convergence rates against Q's baseline. This isolates whether information-sharing alone---without cooperative incentives---addresses the strategic misalignment Q identified.

\section*{Evidence Quality}

The connexion rests on abstract-level analysis of two papers. No empirical evidence supports the transfer hypothesis. The ledger entries are analytical notes, not experimental results. Evidence is thin.

\end{document}
```